\documentclass[11pt]{article}
\usepackage[letterpaper,margin=0.83in,headheight=16pt]{geometry}
\usepackage[T1]{fontenc}
\usepackage[utf8]{inputenc}
\usepackage{lmodern}
\usepackage{microtype}
\usepackage{amsmath,amssymb}
\usepackage{xcolor}
\usepackage{hyperref}
\usepackage{enumitem}
\usepackage{booktabs,longtable,tabularx,array}
\usepackage{fancyhdr}
\usepackage{titlesec}
\usepackage{listings}
\usepackage{tikz}
\usepackage{float}
\definecolor{suiteblue}{HTML}{17365D}
\definecolor{suitelight}{HTML}{EAF1F8}
\definecolor{suitegray}{HTML}{556070}
\hypersetup{colorlinks=true,linkcolor=suiteblue,urlcolor=suiteblue,breaklinks=true}
\titleformat{\section}{\Large\bfseries\color{suiteblue}}{\thesection}{0.55em}{}
\titleformat{\subsection}{\large\bfseries\color{suiteblue}}{\thesubsection}{0.55em}{}
\setlist[itemize]{itemsep=0.25em,topsep=0.35em,leftmargin=1.4em}
\setlist[enumerate]{itemsep=0.25em,topsep=0.35em,leftmargin=1.6em}
\setlength{\parindent}{0pt}
\setlength{\parskip}{0.47em}
\emergencystretch=2em
\lstset{basicstyle=\ttfamily\footnotesize,breaklines=true,frame=single,columns=fullflexible,showstringspaces=false}
\pagestyle{fancy}\fancyhf{}
\lhead{Game Design \& Engineering}\rhead{Chapter Study Guide}\cfoot{\thepage}
\begin{document}
{\huge\bfseries\color{suiteblue} Chapter 29: The Team Builds a Game with Technology}\par
\vspace{0.15in}
{\normalsize\color{suitegray} Primary source: Schell, third edition, Chapter 29, printed p. 497}\par
\vspace{0.15in}\hrule\vspace{0.15in}

\section{Scope and provenance}
Source boundary: "Book-aligned" sections are original, selective study prompts grounded in Jesse Schell, The Art of Game Design: A Book of Lenses, third edition, Chapter 29 (printed page 497). "Engineering extension" and numeric acceptance criteria are supplementary examples, not prescriptions from the book.


\section{Book-aligned study map}
The chapter is read through the following key concepts; the descriptions below are an original learning map rather than a reproduction of the chapter.


\subsection{Foundational versus decorative technology}
Explain foundational versus decorative technology in the context of a concrete player action. Identify its influence on perceived experience, the assumptions required for the concept to apply, and a counterexample in which a different design choice would work better.


\subsection{Technical choice}
Explain technical choice in the context of a concrete player action. Identify its influence on perceived experience, the assumptions required for the concept to apply, and a counterexample in which a different design choice would work better.


\subsection{Future uncertainty}
Explain future uncertainty in the context of a concrete player action. Identify its influence on perceived experience, the assumptions required for the concept to apply, and a counterexample in which a different design choice would work better.


\section{Learning outcomes}
\begin{itemize}
\item Define and distinguish the chapter topics: foundational versus decorative technology, technical choice, future uncertainty.
\item Apply the chapter ideas to one complete playable decision, not merely a visual mockup.
\item Evaluate at least one unintended consequence of the design for a different player population.
\item Explain which claims need direct playtest evidence and which are design assumptions.
\item Translate an experiential claim into an observable or inspectable design artifact.
\end{itemize}

\section{Reading and analysis questions}
\begin{enumerate}
\item Which player-facing problem connects foundational versus decorative technology to future uncertainty?
\item How would you distinguish a design failure from an execution failure when applying foundational versus decorative technology?
\item Which of your conclusions come from the author, and which are an interpretation based on your project?
\item What would a new player do differently from an expert player, and why?
\item What competing design goals would force a trade-off rather than a universal best answer?
\end{enumerate}

\section{Design laboratory -- book-inspired}
Compare two candidate engines or runtime stacks for one design constraint.

\begin{enumerate}
\item Write the player goal and hypothesized experience before implementation.
\item Build the smallest playable or paper prototype that tests the claim.
\item Observe at least three distinct attempts; record outcomes without prompting the player.
\item Revise one variable and document both the intended and actual effect.
\end{enumerate}

\section{Supplementary engineering extension}
Prototype critical path, profile performance, document exit criteria.


\subsection{Reproducible engineering protocol}
\begin{enumerate}
\item Define a stable identifier for the design requirement and record its originating chapter concept.
\item State assumptions, constraints, invariants, failure modes, and instrumentation boundaries.
\item Define one functional test and one player-facing validation question; do not treat them as interchangeable.
\item Record implementation revision, test environment, and evidence before making the next change.
\end{enumerate}

\section{Submission and evidence}
Technology evaluation matrix and architecture decision record.


\section{Review rubric (20 points)}
\begin{tabularx}{\linewidth}{@{}Xcc@{}}\toprule
Criterion & Maximum & Awarded \\ \midrule
Source concept interpreted accurately & 4 & \rule{0.45in}{0.4pt}\\
Problem and audience defined & 4 & \rule{0.45in}{0.4pt}\\
Prototype and validity of evidence & 4 & \rule{0.45in}{0.4pt}\\
Engineering rigor and reproducibility & 4 & \rule{0.45in}{0.4pt}\\
Ethics, accessibility, and reflection & 4 & \rule{0.45in}{0.4pt}\\\bottomrule
\end{tabularx}

\section{Reflection and next dependency}
Record the original design assumption, observation that most challenged it, the resulting design decision, and a question to carry into the following chapter. Re-read relevant lenses only when the corresponding design choice has changed.


\end{document}
